\chapter{Apprentissage par renforcement pour l'exécution et le market making}\label{ch:ml:reinforcement-learning-for-execution-and-market-making}

L'exécution et le market making ont tous deux des \omterm{def:ml:reinforcement-learning-foundations:mdp}{politiques} optimales sous forme fermée dans leurs modèles de manuel~: le calendrier d'Almgren--Chriss
(livre~10, chapitre~14) et les cotations tenant compte de l'inventaire (livre~11, chapitre~3). Ce qui justifie d'apprendre une \omterm{def:ml:reinforcement-learning-foundations:mdp}{politique}
est tout ce que le modèle laisse de côté~: le risque, comme le montre le chapitre~17, est qu'une \omterm{def:ml:reinforcement-learning-foundations:mdp}{politique} apprise exploite ce que son
monde d'entraînement modélise mal. Ce chapitre mesure les deux sur des problèmes dont l'optimum est connu. Un \omterm{def:ml:reinforcement-learning-for-execution-and-market-making:dqn}{réseau Q profond} apprend le
calendrier d'Almgren--Chriss à 0.4~points de base près~; transféré sur un marché dont l'impact est trois fois plus élevé, il perd 9.1
points de base par rapport à l'optimum de ce marché, soit quatre fois ce que perd le calendrier du manuel dans le modèle erroné. Entraîné sur un impact randomisé
et autorisé à voir l'impact de ses propres exécutions, il perd 0.9 là-bas et 0.6 chez lui. Un apprenant tabulaire pour le market
making, avec 1.6~million de pas, reste en deçà de la forme fermée et d'une cotation constante.

\section{Conception de l'état, de l'action et de la récompense}

\begin{definition}[Façonnage de la récompense, masquage d'actions]\label{def:ml:reinforcement-learning-for-execution-and-market-making:design}
Le \emph{façonnage de la récompense}\index{façonnage de la récompense} (reward shaping) modifie les \omterm{def:ml:reinforcement-learning-foundations:mdp}{récompenses} dont un agent apprend sans changer la \omterm{def:ml:reinforcement-learning-foundations:mdp}{politique}
optimale, par exemple en retirant un terme dont l'espérance est connue pour être nulle, ou en ajoutant une différence de potentiel, pour rendre
l'apprentissage plus rapide ou plus stable. Le \emph{masquage d'actions}\index{masquage d'actions} (action masking) retire du choix de l'agent et du
maximum de sa cible d'apprentissage les actions interdites dans un état donné (vendre plus que ce qui est détenu,
coter un côté qui violerait une limite d'inventaire).
\end{definition}

La conception détermine l'essentiel de ce qu'un agent peut apprendre. L'agent d'exécution vend une unité en dix pas, par lots d'un vingtième~;
son état est le temps, les avoirs et un nombre décrit ci-dessous~; son action est le nombre de lots à vendre, toute vente
au-delà des avoirs étant masquée et le reliquat entier étant imposé au dernier pas. Sa \omterm{def:ml:reinforcement-learning-foundations:mdp}{récompense} est l'objectif d'Almgren--Chriss
découpé par pas~: moins le coût d'impact temporaire $\eta u^2/\tau$ de la vente, moins la charge de risque $\lambda\sigma^2\tau x^2$ sur
ce qui reste détenu, avec $\eta = 0.001$, $\sigma = 0.02$ et $\lambda = 10$ (en points de base, le TWAP coûte 21.4 et l'optimum
18.85). Le P\&L réalisé des avoirs sur le pas, $\sigma\sqrt\tau Z x$, a une espérance nulle~; le retirer relève
du \omterm{def:ml:reinforcement-learning-for-execution-and-market-making:design}{façonnage de la récompense}. Le laisser n'est pas anodin~: avec ce bruit dans sa \omterm{def:ml:reinforcement-learning-foundations:mdp}{récompense}, le même réseau, après les mêmes 1\,500
épisodes, a appris à vendre le même montant à chaque pas, c'est-à-dire exactement le TWAP, à 2.55 points de base de l'optimum.

\section{Réseaux Q profonds et leurs stabilisateurs}

\begin{definition}[Réseau Q profond, rejeu d'expérience, réseau cible]\label{def:ml:reinforcement-learning-for-execution-and-market-making:dqn}
Un \emph{réseau Q profond}\index{réseau Q profond} (deep Q-network, DQN) est le \omterm{def:ml:reinforcement-learning-foundations:td}{Q-learning} (chapitre~17) dans lequel la \omterm{def:ml:reinforcement-learning-foundations:bellman}{fonction de valeur d'action} est représentée par
un \omterm{def:ml:neural-networks-for-noisy-tabular-data:mlp}{réseau de neurones}, entraîné par des pas de gradient stochastiques sur l'erreur de différence temporelle (Mnih et coauteurs, 2015). Deux
stabilisateurs le rendent opérant~: le \emph{rejeu d'expérience}\index{rejeu d'expérience} (experience replay) conserve les transitions passées et entraîne sur des lots tirés au hasard
parmi elles, ce qui rompt la corrélation entre pas consécutifs~; un \emph{réseau cible}\index{réseau cible} (target network) est une copie du
réseau, mis à jour seulement toutes les quelques centaines de pas, et servant à calculer les cibles d'apprentissage afin qu'elles ne bougent pas à chaque
mise à jour.
\end{definition}

Le DQN du chapitre (\cref{lst:ml:rlt:dqn}) est un réseau à deux couches de 64 unités à 21 sorties, une mémoire de rejeu de 20\,000
transitions, un \omterm{def:ml:reinforcement-learning-for-execution-and-market-making:dqn}{réseau cible} recopié toutes les 200 mises à jour et une exploration qui décroît de 1 à 0.05 sur la première moitié de
l'entraînement. Après 1\,500 épisodes, sa \omterm{def:ml:reinforcement-learning-foundations:mdp}{politique} gloutonne coûte 19.29 points de base, soit 0.44 au-dessus de l'optimum continu
d'Almgren--Chriss~; 0.20 de cet écart tient à la grille des lots, dont l'optimum propre coûte 19.05. La maximisation sur des estimations bruitées biaise le \omterm{def:ml:reinforcement-learning-foundations:td}{Q-learning}
à la hausse~; le double \omterm{def:ml:reinforcement-learning-foundations:td}{Q-learning}, qui choisit l'action avec un réseau et l'évalue avec l'autre (van Hasselt, Guez et
Silver, 2016), est le remède habituel et le choix de conception qu'ont retenu Ning et coauteurs (2021) pour l'exécution.

\section{Apprendre à exécuter}

\begin{definition}[Randomisation de domaine]\label{def:ml:reinforcement-learning-for-execution-and-market-making:dr}
La \emph{randomisation de domaine}\index{randomisation de domaine} (domain randomisation) entraîne une \omterm{def:ml:reinforcement-learning-foundations:mdp}{politique} sur de nombreuses versions d'un simulateur dont les paramètres incertains
sont tirés au hasard à chaque épisode, afin que la \omterm{def:ml:reinforcement-learning-foundations:mdp}{politique} fonctionne sur l'ensemble de ces versions au lieu d'en exploiter une seule (Tobin et
coauteurs, 2017).
\end{definition}

La faiblesse du manuel tient à ses paramètres~: l'impact est estimé avec erreur et varie avec le marché. Trois mondes testent les
agents~: celui du modèle ($\eta = 0.001$), un marché dont l'impact est trois fois plus élevé, et un autre où il est divisé par trois~; chacun a son propre
optimum d'Almgren--Chriss, et chaque calendrier est évalué exactement par rapport à celui-ci (\cref{fig:ml:rlt:gaps}). La troisième variable d'état
est l'impact que l'agent a observé dans ses propres exécutions jusqu'ici, sous la forme du logarithme de son rapport à la valeur du modèle (nul avant
la première vente)~: un véritable algorithme d'exécution mesure son impact au fil de ses transactions.

\begin{omfigure}
\begin{tikzpicture}
\begin{axis}[width=12cm,height=6cm,ybar,bar width=7pt,ymin=0,ymax=10,ylabel={écart à l'optimum du monde (pb)},
  area legend,xtick={0,1,2},xticklabels={impact du modèle,impact $\times3$,impact $/3$},tick label style={font=\small},
  label style={font=\small},xmin=-0.5,xmax=2.5,grid=major,grid style={gray!25},legend columns=3,legend cell align=left,
  legend style={font=\footnotesize,at={(0.5,-0.18)},anchor=north,draw=none,fill=none},table/col sep=comma]
\addplot[fill=black!60,draw=none,bar shift=-16pt] table[x=i,y=acmodel]{figdata/ml/18-reinforcement-learning-for-execution-and-market-making/gaps.csv};
\addplot[fill=gray!40,draw=none,bar shift=-8pt] table[x=i,y=twap]{figdata/ml/18-reinforcement-learning-for-execution-and-market-making/gaps.csv};
\addplot[fill=omThm,draw=none,bar shift=0pt] table[x=i,y=dqn]{figdata/ml/18-reinforcement-learning-for-execution-and-market-making/gaps.csv};
\addplot[fill=omThm!40,draw=none,bar shift=8pt] table[x=i,y=noisy]{figdata/ml/18-reinforcement-learning-for-execution-and-market-making/gaps.csv};
\addplot[fill=omDef,draw=none,bar shift=16pt] table[x=i,y=dr]{figdata/ml/18-reinforcement-learning-for-execution-and-market-making/gaps.csv};
\legend{Almgren--Chriss (modèle),TWAP,DQN,DQN à récompense bruitée,DQN randomisé}
\end{axis}
\end{tikzpicture}

\omcaption{Écart de l'objectif de chaque calendrier à l'optimum d'Almgren--Chriss du monde dans lequel il opère (points de base). La forme fermée
et le premier DQN ont été construits pour l'impact du modèle~; le DQN randomisé a été entraîné sur un impact tiré sur un facteur
trois dans un sens comme dans l'autre. Données~: \texttt{ml\_rltrade.execution}.}\label{fig:ml:rlt:gaps}
\end{omfigure}

Le DQN entraîné sur le modèle est l'avertissement du chapitre. Dans son propre monde, il est à 0.44 de l'optimum~; là où l'impact est trois fois
plus élevé, il en est à 9.09, et là où il est plus faible, à 4.60~: il n'a jamais vu sa troisième entrée prendre ces valeurs, et un réseau
hors de sa distribution d'entraînement fait n'importe quoi. La forme fermée construite sur le même impact erroné perd 2.20 et 1.21~:
une formule se dégrade en douceur là où un réseau peut ne pas le faire. Le DQN randomisé, entraîné pendant 3\,000 épisodes avec un impact tiré
à nouveau à chaque épisode, a appris à lire ses exécutions~: il vend d'abord un peu plus lentement que l'optimum du modèle, puis
suit l'impact qu'il observe, et ses écarts sont de 0.55, 0.86 et 1.24~: le meilleur de tous dans le monde à fort impact, et à moins de 0.03 du meilleur (la
forme fermée, par hasard du sens de son erreur) dans le monde à faible impact.
La \Cref{fig:ml:rlt:schedules} montre les calendriers dans le monde à fort impact.

\begin{omfigure}
\begin{tikzpicture}
\begin{axis}[width=11cm,height=6cm,xlabel={pas},ylabel={avoirs (part de l'ordre)},tick label style={font=\small},
  label style={font=\small},xmin=0,xmax=10,ymin=0,ymax=1.02,grid=major,grid style={gray!25},legend columns=3,legend cell align=left,
  legend style={font=\footnotesize,at={(0.5,-0.25)},anchor=north,draw=none,fill=none,
  /tikz/every even column/.append style={column sep=6pt}},table/col sep=comma]
\addplot[black,ultra thick] table[x=step,y=optimum]{figdata/ml/18-reinforcement-learning-for-execution-and-market-making/schedules.csv};
\addplot[black!50,thick,dashed] table[x=step,y=acmodel]{figdata/ml/18-reinforcement-learning-for-execution-and-market-making/schedules.csv};
\addplot[gray,thick,dotted] table[x=step,y=twap]{figdata/ml/18-reinforcement-learning-for-execution-and-market-making/schedules.csv};
\addplot[omThm,very thick,mark=*,mark size=1.4pt] table[x=step,y=dqn]{figdata/ml/18-reinforcement-learning-for-execution-and-market-making/schedules.csv};
\addplot[omDef,very thick,mark=square*,mark size=1.4pt] table[x=step,y=dr]{figdata/ml/18-reinforcement-learning-for-execution-and-market-making/schedules.csv};
\legend{optimum pour impact $\times3$,Almgren--Chriss (modèle),TWAP,DQN,DQN randomisé}
\end{axis}
\end{tikzpicture}

\omcaption{Avoirs sur les dix pas dans le monde à impact trois fois celui du modèle. Données~:
\texttt{ml\_rltrade.agents}.}\label{fig:ml:rlt:schedules}
\end{omfigure}

\section{Apprendre à faire du market making}

Le market making dans le modèle du livre~11 (chapitre~3) admet une solution exacte~: un prix milieu qui diffuse, des exécutions à la profondeur $\delta$
avec une intensité $Ae^{-k\delta}$, une pénalité courante $\phi q^2$ sur l'inventaire et une pénalité terminale $\alpha q^2$, ici avec
$A = 140$, $k = 1.5$, $\sigma = 2$, $\phi = 0.02$, $\alpha = 0.01$ et un inventaire compris dans $\pm10$. Les profondeurs optimales au départ
sont de 0.68 des deux côtés lorsque la position est nulle et de 0.57 et 0.79 lorsque l'on est long de cinq unités (plus serrée du côté qui réduit l'inventaire). Le
dispositif d'apprentissage est un \omterm{def:ml:reinforcement-learning-foundations:td}{Q-learning} tabulaire (\cref{lst:ml:rlt:mm}) sur 20 intervalles de temps et les 21 niveaux d'inventaire, qui choisit la profondeur de chaque côté
parmi six valeurs, le P\&L d'inventaire de moyenne nulle étant retiré de la \omterm{def:ml:reinforcement-learning-foundations:mdp}{récompense} par façonnage et les côtés qui franchiraient la limite d'inventaire étant masqués.
Le \Cref{tab:ml:rlt:mm} évalue chaque \omterm{def:ml:reinforcement-learning-foundations:mdp}{politique} avec le simulateur de market making de l'entreprise sur 5\,000 trajectoires communes.

\begin{table}[htbp]
\centering\small
\begin{tabular}{@{}lccccc@{}}
\toprule
\omterm{def:ml:reinforcement-learning-foundations:mdp}{politique} & objectif & (e.t.) & écart-type du P\&L & $|q_T|$ & exécutions\\
\midrule
optimum exact (livre~11) & 67.95 & 0.12 & 8.53 & 2.67 & 101.8\\
Avellaneda--Stoikov, $\gamma = 0.1$ & 64.73 & 0.09 & 6.50 & 2.26 & 97.0\\
profondeur constante $1/k$ & 66.67 & 0.17 & 11.61 & 5.03 & 101.0\\
\omterm{def:ml:reinforcement-learning-foundations:td}{Q-learning}, 2\,000 épisodes & 62.34 & 0.12 & 8.31 & 4.40 & 91.7\\
\omterm{def:ml:reinforcement-learning-foundations:td}{Q-learning}, 8\,000 épisodes & 62.84 & 0.12 & 8.53 & 4.38 & 93.1\\
\bottomrule
\end{tabular}
\caption{Market making pour une unité de temps~: moyenne du P\&L diminuée des pénalités d'inventaire courante et terminale (avec son
erreur type), écart-type du P\&L, inventaire moyen de clôture et exécutions par trajectoire, sur 5\,000 trajectoires à
nombres aléatoires communs. Données~: \texttt{ml\_rltrade.market\_making}.}\label{tab:ml:rlt:mm}
\end{table}

Le dispositif d'apprentissage progresse lentement et reste sous la cotation constante~: 1.6 million de pas d'expérience ne suffisent pas à classer six
profondeurs dont les espérances diffèrent de moins que le bruit d'une exécution. La fourchette gagnée à chaque pas est un petit signal sous un
bruit important, c'est là le problème du teneur de marché en général. Les formules fermées ou les \omterm{def:ml:reinforcement-learning-foundations:mdp}{politiques} paramétriques ajustées par recherche de \omterm{def:ml:reinforcement-learning-foundations:mdp}{politique} dans
un simulateur (une profondeur qui est une droite en fonction de l'inventaire a deux paramètres, là où le tableau compte 420 états et 36 actions) sont ce qui fonctionne à cette échelle~;
l'apprentissage ne gagne sa place que sur les parties que le modèle laisse de côté, comme la position dans la file, la sélection adverse et plusieurs plateformes de négociation,
là où aucune formule fermée n'existe.

\section{Ce qui fonctionne dans les travaux publiés}

Les résultats sont plus minces que l'attention qu'on leur porte. Nevmyvaka, Feng et Kearns (2006) ont appliqué l'\omterm{def:ml:reinforcement-learning-foundations:mdp}{apprentissage par renforcement} à l'exécution sur
un an et demi de données d'ordres à cours limité à la milliseconde issues du NASDAQ, avec un espace d'états factorisé pour que l'apprentissage reste traitable~;
Ning et ses coauteurs (2021) ont entraîné un double \omterm{def:ml:reinforcement-learning-for-execution-and-market-making:dqn}{réseau Q profond} avec \omterm{def:ml:reinforcement-learning-for-execution-and-market-making:dqn}{rejeu d'expérience} sur les variables du carnet d'ordres pour neuf actions, et
ils ont constaté qu'il battait l'indice de référence standard sur la plupart d'entre elles~; Spooner et ses coauteurs
(2018) ont construit un simulateur de carnet d'ordres à cours limité très fidèle et un teneur de marché à différences temporelles doté d'une \omterm{def:ml:reinforcement-learning-foundations:mdp}{récompense} sur mesure qui
contrôle le risque d'inventaire, et qui a battu des indices de référence simples ainsi qu'une approche d'\omterm{def:ml:online-learning-and-drift:rls}{apprentissage en ligne} dans ce simulateur. Hambly, Xu et Yang
(2023) dressent un panorama du domaine. Le fil conducteur est le simulateur~: les résultats sont rapportés dans un simulateur, souvent celui des auteurs, et
le passage au trading en conditions réelles est rarement documenté publiquement. Un desk doit les lire comme une indication que les méthodes fonctionnent là où
le simulateur est fidèle, et construire ses propres éléments probants avec les tests du chapitre~17 et de ce chapitre~: indices de référence exacts, plusieurs
mondes différents, \omterm{def:ml:reinforcement-learning-foundations:ope}{évaluation hors politique} (off-policy evaluation) sur ses propres journaux.

\begin{method}[Déployer une politique d'exécution ou de cotation apprise]\label{met:ml:rlt:method}
\begin{enumerate}
\item Comparer à la formule fermée du modèle de manuel, dans ce même modèle~: un apprenant qui ne parvient pas à l'égaler là n'est pas
prêt.
\item Façonner la \omterm{def:ml:reinforcement-learning-foundations:mdp}{récompense} pour retirer le bruit connu de moyenne nulle, masquer les actions interdites et donner à l'état les signaux dont un
trader se servirait (exécutions, impact observé, position dans la file).
\item Entraîner sur des simulateurs randomisés dont les paramètres couvrent leur erreur d'estimation, et évaluer sur des jeux de paramètres
réservés, jamais seulement dans celui de l'entraînement.
\item Borner la \omterm{def:ml:reinforcement-learning-foundations:mdp}{politique} (limites de participation, limites d'inventaire, repli sur la formule fermée) et comparer en conditions réelles avec
le dispositif en place, par \omterm{def:ml:reinforcement-learning-foundations:ope}{évaluation hors politique} et par de petits essais contrôlés.
\end{enumerate}
\end{method}

\section{Tutoriel~: battre une forme fermée}\label{tut:ml:reinforcement-learning-for-execution-and-market-making:tut}

\begin{tutorial}
\textbf{Objectif.} Entraîner un DQN sur la liquidation d'Almgren--Chriss avec et sans \omterm{def:ml:reinforcement-learning-for-execution-and-market-making:dr}{randomisation de domaine}, et l'évaluer dans trois
mondes~; entraîner un teneur de marché tabulaire et l'évaluer par rapport à l'optimum exact. \textbf{État final~:}
les \cref{fig:ml:rlt:gaps,fig:ml:rlt:schedules}, le \cref{tab:ml:rlt:mm}.
\begin{tutsteps}
\item \textbf{Une mise à jour DQN avec rejeu, \omterm{def:ml:reinforcement-learning-for-execution-and-market-making:dqn}{réseau cible} et masquage.}
\omcode{code/firm/rltrade/firm_rltrade.py}{113}{143}{La boucle d'entraînement du Q-learning profond.}\label{lst:ml:rlt:dqn}
\item \textbf{\omterm{def:ml:reinforcement-learning-foundations:td}{Q-learning} pour le market making.}
\omcode{code/firm/rltrade/firm_rltrade.py}{228}{243}{Q-learning tabulaire à pas décroissants.}\label{lst:ml:rlt:mm}
\item \textbf{Lancer} \texttt{ml\_rltrade.execution()} (environ une minute sur un cœur), \texttt{market\_making()} et
\texttt{fig\_rltrade.py}.
\end{tutsteps}
\textbf{Pour aller plus loin.} Utiliser le double \omterm{def:ml:reinforcement-learning-foundations:td}{Q-learning}~; remplacer la table du teneur de marché par une \omterm{def:ml:reinforcement-learning-foundations:mdp}{politique} à deux paramètres ajustée par
recherche de \omterm{def:ml:reinforcement-learning-foundations:mdp}{politique}~; ajouter une action d'ordre passif dont les exécutions proviennent de la position dans la file du simulateur de bourse du livre~10.
\end{tutorial}

\section{Construction~: environnements de trading et apprenants}\label{bld:ml:reinforcement-learning-for-execution-and-market-making:rltrade}

\begin{build}
\bfield{Rôle} \omterm{def:ml:reinforcement-learning-foundations:mdp}{Politiques} d'exécution et de market making apprises là où leur optimum est connu, et testées loin de cet optimum.
\bfield{Interface} \texttt{ACEnv(X, T, n, lots, eta, sigma, lam, eta\_range, noise)} avec \texttt{reset}, \texttt{step},
\texttt{mask}~; \texttt{DQN}, \texttt{train\_dqn(env, episodes, seed)}, \texttt{dqn\_schedule(env, net, eta)},
\texttt{grid\_optimum}, \texttt{ac\_objective} (sur \texttt{firm.acexec})~; \texttt{MMEnv}, \texttt{q\_learning\_mm},
\texttt{mm\_policy}, \texttt{mm\_objective} (avec \texttt{firm.invmm.simulate}).
\bfield{Règles} Toute \omterm{def:ml:reinforcement-learning-foundations:mdp}{politique} apprise est évaluée exactement ou sur nombres aléatoires communs par rapport à la formule fermée du monde dans lequel
elle est testée~; les paramètres d'entraînement et ceux de test sont présentés séparément.
\bfield{Tests d'acceptation} \texttt{code/firm/rltrade/tests/}~: les \omterm{def:ml:reinforcement-learning-foundations:mdp}{récompenses} de l'environnement totalisent l'objectif d'Almgren--Chriss de
l'échéancier suivi~; le masquage n'autorise jamais une vente illégale~; l'optimum de la grille de lots n'est pas moins bon que tout échéancier énuméré~;
l'entraînement DQN est déterministe et bat le TWAP dans le monde du modèle~; les variations d'inventaire de l'environnement de market making coïncident
avec les exécutions, de probabilité $Ae^{-k\delta}\,dt$ par côté.
\bfield{Pour aller plus loin} DQN double~; recherche de \omterm{def:ml:reinforcement-learning-foundations:mdp}{politique} pour la cotation~; l'enveloppe du simulateur de bourse (livre~10, chapitre~26) avec la
position dans la file comme état.
\end{build}

\begin{omsources}
\item V.~Mnih et ses coauteurs, «~Human-level control through deep reinforcement learning~», \emph{Nature} 518, 2015.
\item H.~van Hasselt, A.~Guez et D.~Silver, «~Deep reinforcement learning with double Q-learning~», arXiv:1509.06461, 2016.
\item J.~Tobin et ses coauteurs, «~Domain randomization for transferring deep neural networks from simulation to the real
world~», arXiv:1703.06907, 2017.
\item Y.~Nevmyvaka, Y.~Feng et M.~Kearns, «~Reinforcement learning for optimized trade execution~», ICML, 2006.
\item B.~Ning et ses coauteurs, «~Double deep Q-learning for optimal execution~», \emph{Applied Mathematical Finance}, 2021.
\item T.~Spooner, J.~Fearnley, R.~Savani et A.~Koukorinis, «~Market making via reinforcement learning~», arXiv:1804.04216,
2018.
\item B.~Hambly, R.~Xu et H.~Yang, «~Recent advances in reinforcement learning in finance~», arXiv:2112.04553, 2023.
\item R.~Almgren et N.~Chriss, «~Optimal execution of portfolio transactions~», \emph{Journal of Risk} 3(2), 2001.
\item M.~Avellaneda et S.~Stoikov, «~High-frequency trading in a limit order book~», \emph{Quantitative Finance} 8(3), 2008.
\end{omsources}

\section{Exercices}

\begin{exercise}[$\star$]\label{exo:ml:reinforcement-learning-for-execution-and-market-making:1}
Calculer l'objectif du TWAP dans le monde modèle~: dix ventes de 0.1 avec $\eta = 0.001$, $\tau = 0.1$, et la pénalité de risque avec
$\sigma = 0.02$, $\lambda = 10$. Le donner en points de base.
\end{exercise}

\begin{exercise}[$\star$]\label{exo:ml:reinforcement-learning-for-execution-and-market-making:2}
À quelle profondeur $\delta Ae^{-k\delta}$, la fourchette espérée gagnée par unité de temps d'un côté, atteint-elle son maximum~? Comparer avec la
profondeur plate de l'optimum exact.
\end{exercise}

\begin{exercise}[$\star$]\label{exo:ml:reinforcement-learning-for-execution-and-market-making:3}
Pourquoi le masque doit-il aussi intervenir dans la cible $\max_{u'}Q(x', u')$, et pas seulement dans le choix de l'action~?
\end{exercise}

\begin{exercise}[$\star\star$]\label{exo:ml:reinforcement-learning-for-execution-and-market-making:4}
Montrer que retirer $\sigma\sqrt\tau Zx$ de la \omterm{def:ml:reinforcement-learning-foundations:mdp}{récompense} de chaque pas laisse la \omterm{def:ml:reinforcement-learning-foundations:mdp}{politique} optimale inchangée. Pourquoi l'agent bruité
a-t-il fini au TWAP~?
\end{exercise}

\begin{exercise}[$\star\star$]\label{exo:ml:reinforcement-learning-for-execution-and-market-making:5}
Pourquoi le DQN entraîné sur le modèle a-t-il fait moins bien que la mauvaise forme fermée dans le monde à fort impact~?
\end{exercise}

\begin{exercise}[$\star\star$]\label{exo:ml:reinforcement-learning-for-execution-and-market-making:6}
\emph{Trouver la faille.} «~Nous avons entraîné le teneur de marché avec une \omterm{def:ml:reinforcement-learning-for-execution-and-market-making:dr}{randomisation de domaine} sur la volatilité, et il a été rentable dans chacune des
10\,000 simulations randomisées du simulateur de bourse, donc il est robuste.~»
\end{exercise}

\begin{exercise}[$\star\star\star$]\label{exo:ml:reinforcement-learning-for-execution-and-market-making:7}
\emph{Programmation.} Remplacer la table du teneur de marché par la \omterm{def:ml:reinforcement-learning-foundations:mdp}{politique} $\delta_{b,a} = c\pm dq$ et choisir $(c, d)$ sur une grille selon
l'objectif moyen du simulateur sur 2\,000 trajectoires. À quelle distance de l'optimum exact parvient-on~?
\end{exercise}

\begin{exercise}[$\star\star\star$]\label{exo:ml:reinforcement-learning-for-execution-and-market-making:8}
Établir le calendrier discret d'Almgren--Chriss pour l'objectif du chapitre en écrivant les conditions du premier ordre sur les
positions $x_1,\dots,x_{n-1}$.
\end{exercise}

\section{Problème~: Battre une forme fermée}

\begin{problem}[{Problème du week-end --- quand apprendre une politique}]\label{pb:ml:reinforcement-learning-for-execution-and-market-making:1}
Les problèmes de liquidation et de teneur de marché du chapitre.

\medskip\noindent\textbf{Partie I --- Conception.}
\begin{enumerate}
\item Quels sont l'état, les actions et la \omterm{def:ml:reinforcement-learning-foundations:mdp}{récompense} de l'agent d'exécution~?
\item Pourquoi le bruit du P\&L est-il retiré de la \omterm{def:ml:reinforcement-learning-foundations:mdp}{récompense}, et que s'est-il passé lorsqu'il ne l'était pas~?
\item Qu'est-ce que le \omterm{def:ml:reinforcement-learning-for-execution-and-market-making:design}{masquage d'actions} empêche~?
\item Pourquoi donner à l'agent l'impact qu'il a observé~?
\end{enumerate}

\medskip\noindent\textbf{Partie II --- Exécution.}
\begin{enumerate}[resume]
\item Que coûtent le TWAP, l'optimum de la grille de lots et le DQN dans le monde modèle~?
\item Quels sont les écarts de chaque calendrier dans les trois mondes~?
\item Pourquoi l'agent randomisé perd-il un peu chez lui~?
\item Que changerait un prix qui tend pendant l'ordre~?
\end{enumerate}

\medskip\noindent\textbf{Partie III --- Market making.}
\begin{enumerate}[resume]
\item Quelles sont les profondeurs de l'optimum exact au départ~?
\item Comment les cinq \omterm{def:ml:reinforcement-learning-foundations:mdp}{politiques} se classent-elles~?
\item Pourquoi l'apprenant tabulaire reste-t-il en dessous d'une cotation constante~?
\item Qu'apprendrait-on, et que garderait-on sous forme fermée~?
\end{enumerate}

\medskip\noindent\textbf{Partie IV --- Le verdict.}
\begin{enumerate}[resume]
\item Énoncer le \emph{résultat nommé}~: l'écart de la \omterm{def:ml:reinforcement-learning-foundations:mdp}{politique} apprise à l'optimum d'Almgren--Chriss dans le monde modèle et
dans les mondes mal spécifiés, avec et sans \omterm{def:ml:reinforcement-learning-for-execution-and-market-making:dr}{randomisation de domaine}.
\item Que montrent les données publiées, et que ne montrent-elles pas~?
\item Comment déploierait-on l'agent d'exécution randomisé~?
\item Quel est le repli si l'impact observé par l'agent sort de sa plage d'entraînement~?
\item Pourquoi une table ne passe-t-elle pas à l'échelle d'un vrai carnet d'ordres~?
\item Quel stabilisateur compte le plus dans le DQN ici, et comment le vérifier~?
\item À quoi doit ressembler la comparaison avec l'algorithme en place~?
\item En une phrase~: quand une \omterm{def:ml:reinforcement-learning-foundations:mdp}{politique} apprise vaut-elle plus qu'une forme fermée~?
\end{enumerate}
\end{problem}

\section{Questions d'entretien}

\begin{interviewq}[$\star$ \iqroles{ingénieur ML}]\label{iq:ml:reinforcement-learning-for-execution-and-market-making:1}
À quoi servent le \omterm{def:ml:reinforcement-learning-for-execution-and-market-making:dqn}{rejeu d'expérience} et les réseaux cibles~?
\end{interviewq}

\begin{interviewq}[$\star\star$ \iqroles{chercheur, trader}]\label{iq:ml:reinforcement-learning-for-execution-and-market-making:2}
Comment formuleriez-vous l'exécution optimale comme un problème d'\omterm{def:ml:reinforcement-learning-foundations:mdp}{apprentissage par renforcement}~? Que met-on dans l'état~?
\end{interviewq}

\begin{interviewq}[$\star\star$ \iqroles{chercheur}]\label{iq:ml:reinforcement-learning-for-execution-and-market-making:3}
Qu'est-ce que le \omterm{def:ml:reinforcement-learning-for-execution-and-market-making:design}{façonnage de la récompense}, et comment peut-il mal tourner~?
\end{interviewq}

\begin{interviewq}[$\star\star$ \iqroles{ingénieur ML}]\label{iq:ml:reinforcement-learning-for-execution-and-market-making:4}
Les valeurs Q de votre DQN ne cessent d'augmenter pendant l'entraînement. Que se passe-t-il, et que faites-vous~?
\end{interviewq}

\begin{interviewq}[$\star\star$ \iqroles{chercheur, trader}]\label{iq:ml:reinforcement-learning-for-execution-and-market-making:5}
Utiliseriez-vous l'\omterm{def:ml:reinforcement-learning-foundations:mdp}{apprentissage par renforcement} pour le market making~? Où aiderait-il, et où n'aiderait-il pas~?
\end{interviewq}

\begin{interviewq}[$\star\star\star$ \iqroles{chercheur}]\label{iq:ml:reinforcement-learning-for-execution-and-market-making:6}
Qu'est-ce que la \omterm{def:ml:reinforcement-learning-for-execution-and-market-making:dr}{randomisation de domaine}, et pourquoi peut-elle aussi bien aider que nuire à une \omterm{def:ml:reinforcement-learning-foundations:mdp}{politique} de trading~?
\end{interviewq}
